%=====================================================
\section{Tier 1.3: Mathematical formulation}
\label{sec:tier1_3}

This section describes the governing equations corresponding to the input file
\texttt{tier1\_3\_mode\_I\_load\_unload.i} and the analytical reference
implemented in \texttt{verify\_tier1\_3.py}.

%=====================================================
\subsection{Test purpose}

Tier 1.3 verifies three properties of \texttt{BiLinearMixedModeTraction}
that Tier 1.1 cannot exercise because Tier 1.1 is monotonic:

\begin{enumerate}
\item correctness of the \emph{secant unload/reload} branch
      $t_n = K(1-d)\,\delta_n$ with frozen damage,
\item irreversibility of the damage internal variable
      $\delta_n^{\max}$ (no healing during unload/reload),
\item recovery of the bilinear envelope when reloading past the previously
      attained maximum jump.
\end{enumerate}

%=====================================================
\subsection{Domain, bulk, and interface}

Geometry, mesh discretisation, bulk linear elasticity, and interface
kinematics are identical to Tier 1.1
(see Section~\ref{sec:tier1_1}). In short:

\begin{equation}
\Omega = \Omega_1\cup\Omega_2,\qquad
\Gamma_c = \overline{\Omega_1}\cap\overline{\Omega_2},\qquad
\jump{\bu} = \bu^{+}-\bu^{-},
\end{equation}
\begin{equation}
\delta_n = \jump{\bu}\cdot\bn = \jump{u_y},\qquad
\delta_t = \jump{\bu}\cdot\btan = \jump{u_x},
\end{equation}
\begin{equation}
\nabla\!\cdot\bsigma = \bm{0}\ \ \text{in }\Omega_i,\qquad
\bsigma = \mathbb{C}:\bm{\varepsilon},\qquad
\bsigma\,\bn = \btrac\quad\text{on }\Gamma_c.
\end{equation}

%=====================================================
\subsection{Bilinear law with damage as a state variable}

Damage in the Camanho--D{\'a}vila bilinear model is driven by the
\emph{maximum} jump ever attained at a quadrature point, an internal
variable
\begin{equation}
\delta_n^{\max}(t) \;=\; \max_{\tau\in[0,t]}\;\delta_n(\tau)\,.
\end{equation}
This is what makes the law irreversible: $\delta_n^{\max}$ can only grow.
The damage variable evaluated at $\delta_n^{\max}$ is
\begin{equation}
d\bigl(\delta_n^{\max}\bigr)
\;=\;
\begin{cases}
0,
& \delta_n^{\max} \le \delta_n^{0}, \\[6pt]
\dfrac{\delta_n^{f}\,\bigl(\delta_n^{\max}-\delta_n^{0}\bigr)}
       {\delta_n^{\max}\,\bigl(\delta_n^{f}-\delta_n^{0}\bigr)},
& \delta_n^{0} < \delta_n^{\max} < \delta_n^{f}, \\[10pt]
1,
& \delta_n^{\max} \ge \delta_n^{f},
\end{cases}
\qquad d\in[0,1].
\end{equation}
The traction at the current jump $\delta_n$ uses the \emph{frozen} damage
$d(\delta_n^{\max})$:
\begin{equation}
t_n\bigl(\delta_n,\;\delta_n^{\max}\bigr)
\;=\;
\begin{cases}
\bigl(1-d(\delta_n^{\max})\bigr)\,K\,\delta_n, & \delta_n \ge 0, \\[4pt]
K\,\delta_n, & \delta_n < 0\ \text{(closed contact)}.
\end{cases}
\end{equation}

This piecewise rule contains every regime of interest:
\begin{itemize}
\item \textbf{Loading on the envelope} $\delta_n=\delta_n^{\max}$:
      reduces to the Tier 1.1 bilinear curve.
\item \textbf{Unloading} $\delta_n<\delta_n^{\max}$:
      $\delta_n^{\max}$ is frozen, so the traction is a straight line
      through the origin of slope $K(1-d(\delta_n^{\max}))$.
\item \textbf{Reloading} $\delta_n<\delta_n^{\max}$:
      retraces the same secant.
\item \textbf{Reloading past the previous max} $\delta_n>\delta_n^{\max}$:
      $\delta_n^{\max}$ updates to the new value and the trajectory rejoins
      the bilinear envelope.
\end{itemize}

%=====================================================
\subsection{Numerical parameters}

Same as Tier 1.1:
\begin{equation}
K = 5.0\times10^{3},\qquad
N = 50,\qquad
G_{Ic} = 0.5,
\end{equation}
\begin{equation}
\delta_n^{0} = N/K = 1.0\times10^{-2},\qquad
\delta_n^{f} = 2 G_{Ic}/N = 2.0\times10^{-2}.
\end{equation}

The pre-selected unload point is the \emph{midpoint} of the softening
branch:
\begin{equation}
\delta_n^{\star} \;=\; \tfrac{1}{2}\bigl(\delta_n^{0}+\delta_n^{f}\bigr)
\;=\; 1.5\times10^{-2}.
\end{equation}
Closed-form landmarks at this point:
\begin{equation}
d^{\star} \;=\; d(\delta_n^{\star})
\;=\; \frac{\delta_n^{f}\bigl(\delta_n^{\star}-\delta_n^{0}\bigr)}
           {\delta_n^{\star}\bigl(\delta_n^{f}-\delta_n^{0}\bigr)}
\;=\; \frac{2}{3},
\end{equation}
\begin{equation}
K_{\text{secant}}^{\star} \;=\; K\bigl(1-d^{\star}\bigr)
\;=\; \tfrac{5\,000}{3} \approx 1\,666.67,
\end{equation}
\begin{equation}
t_n^{\star} \;=\; N\,\frac{\delta_n^{f}-\delta_n^{\star}}
                          {\delta_n^{f}-\delta_n^{0}}
\;=\; 25.
\end{equation}

%=====================================================
\subsection{Loading and boundary conditions}

The bottom face is fixed, the top face is constrained in $x$ and pulled in
$y$ by a five-segment piecewise-linear function

\begin{equation}
u_*(t) \;=\;
\begin{cases}
\dfrac{0.010}{0.10}\,t,
& 0.00 \le t \le 0.10,
& \text{elastic loading} \\[6pt]
0.010 + \dfrac{0.005}{0.10}(t-0.10),
& 0.10 < t \le 0.20,
& \text{softening to }\delta_n^{\star} \\[6pt]
0.015 - \dfrac{0.015}{0.20}(t-0.20),
& 0.20 < t \le 0.40,
& \text{unload to zero} \\[6pt]
0.000 + \dfrac{0.015}{0.20}(t-0.40),
& 0.40 < t \le 0.60,
& \text{reload to }\delta_n^{\star} \\[6pt]
0.015 + \dfrac{0.015}{0.40}(t-0.60),
& 0.60 < t \le 1.00,
& \text{continue past }\delta_n^{f}.
\end{cases}
\end{equation}

The end value $u_*(1)=0.030 = 1.5\,\delta_n^{f}$ guarantees full decohesion
($d=1$) by the end of the simulation. The time step is
$\Delta t = 0.005$, giving $200$ samples evenly distributed across the
five branches.

%=====================================================
\subsection{Weak form}

Identical to Tier 1.1:
\begin{equation}
\int_{\Omega_1\cup\Omega_2} \bsigma(\bu):\nabla\bv\,dV
\;+\;
\int_{\Gamma_c} \btrac\bigl(\jump{\bu},\,\delta_n^{\max}\bigr)\cdot
                \jump{\bv}\,dS
\;=\; 0,
\end{equation}
the only difference being that $\btrac$ now depends on the history variable
$\delta_n^{\max}$ stored as a stateful material property in
\texttt{BiLinearMixedModeTraction}.

%=====================================================
\subsection{Verification criterion}

Let $\{(\delta_n^{(k)},\,t_n^{(k)})\}$ be the simulated samples and
construct the running history variable
\begin{equation}
\delta_n^{\max,(k)} \;=\; \max_{j\le k}\;\max\bigl(\delta_n^{(j)},\,0\bigr).
\end{equation}
The closed-form reference is
\begin{equation}
t_n^{\text{ref},(k)}
\;=\;
\bigl(1 - d(\delta_n^{\max,(k)})\bigr)\,K\,\delta_n^{(k)}.
\end{equation}
The test passes if
\begin{equation}
\max_{k}\;\bigl|\,t_n^{(k)} - t_n^{\text{ref},(k)}\,\bigr|
\;<\;10^{-7},
\qquad
\max_{k}\;|t_t^{(k)}|
\;<\;10^{-10},
\end{equation}
and the running history variable is monotone,
\begin{equation}
\delta_n^{\max,(k+1)} \;\ge\; \delta_n^{\max,(k)} \quad\forall\,k.
\end{equation}
The last condition is automatically satisfied by the construction of
$\delta_n^{\max}$, but it is asserted as a sanity check that the
postprocessor output is consistent with damage irreversibility in the
material.

For the parameters above the simulation produced
\begin{equation}
\max_{k}\;\bigl|t_n - t_n^{\text{ref}}\bigr| \;=\; 2.4\times 10^{-12},\qquad
\max_{k}\;|t_t| \;=\; 1.4\times 10^{-14},
\end{equation}
i.e.\ the bilinear law and its irreversibility logic are reproduced to
roughly $13$ significant figures relative to the peak traction $N=50$.

\end{document}
